% Methods: Quantum VQC description — auto-generated by report/paper_sections.py
\subsection{Variational Quantum Circuit (VQC) Feed-Forward Block}
\label{sec:vqc}

The VQC replaces the classical feed-forward sublayer in the second Transformer
encoder layer.
For an input embedding $\mathbf{h} \in \mathbb{R}^{d_\text{model}}$, a linear
projection $W_\text{in} \in \mathbb{R}^{n_q \times d_\text{model}}$ maps to
$n_q = 6$ qubit inputs, scaled to $[-\pi, \pi]$ by a $\tanh(\cdot) \times \pi$
activation.

\paragraph{Hardware-Efficient Ansatz (HEA).}
Each of the $L = 3$ variational layers applies:
\begin{enumerate}
  \item \emph{Data re-uploading}: $\text{RY}(x_i)$ applied to qubit $i$, encoding
        the input at every layer (not only the first), following the expressibility
        analysis of Schuld et al.~\cite{schuld2021effect}.
  \item \emph{Parametric rotations}: $\text{RY}(\theta^{(l)}_{i,1})$ and
        $\text{RZ}(\theta^{(l)}_{i,2})$ per qubit $i$ in layer $l$, totalling
        $2 \times n_q \times L = 36$ variational parameters.
  \item \emph{Ring CNOT entanglement}: $\text{CNOT}(i, (i{+}1) \bmod n_q)$
        for $i = 0,\ldots,n_q{-}1$.
\end{enumerate}
The circuit is implemented in PennyLane~\cite{bergholm2018pennylane} with
\texttt{default.qubit} (classical simulation) and \texttt{diff\_method="backprop"}
for seamless PyTorch gradient flow.
Expectation values $\langle Z_i \rangle$ on all $n_q$ qubits form the circuit
output $\mathbf{q} \in [-1, 1]^{n_q}$, which is mapped back to
$d_\text{model}$ dimensions via a linear layer $W_\text{out}$.

\paragraph{Expressibility and Entanglement.}
We quantify VQC capacity using the Meyer--Wallach measure~\cite{meyerwallach2002}:
expressibility $\mathcal{E} = 1 - \overline{\mathcal{F}}$ (mean pairwise state
fidelity deviation from Haar-random), and entanglement capability
$Q_\text{cap}$ averaged over 2000 random parameter samples.
Full values are reported in Table~\ref{tab:quantum_circuit}.
